\chapter{Requirements Traceability Matrix}
\label{app:traceability}

This appendix links the initial requirements to the implemented artifact and to
the evaluation evidence. It is included to make the PFE contribution easier to
audit: each requirement should be visible in the architecture, implementation,
or evaluation chapters, not only listed at the beginning of the report.

\section{Functional Requirement Traceability}

\begin{longtable}{>{\raggedright\arraybackslash\bfseries}p{1.7cm}
                  >{\raggedright\arraybackslash}p{4.1cm}
                  >{\raggedright\arraybackslash}p{4.0cm}
                  >{\raggedright\arraybackslash}p{4.2cm}}
\caption{Functional requirements mapped to implementation and evaluation.}
\label{tab:functional-traceability}\\
\toprule
\rowcolor{accent!8}
ID & Requirement & Implemented through & Evidence of validation \\
\midrule
\endfirsthead
\toprule
\rowcolor{accent!8}
ID & Requirement & Implemented through & Evidence of validation \\
\midrule
\endhead
\bottomrule
\endfoot
FR-1 & Accept one target URL or a dataset batch. & Operator console, FastAPI run
creation, batch execution records. & Evaluation context table and regression
batch protocol. \\
FR-2 & Classify the starting page. & Classification agent and orchestrator
route graph. & Page-type routing KPI and classification case study. \\
FR-3 & Extract landing-page candidates. & Landing agent, landing inspect tool,
checked-section behavior. & Host-candidate recall KPI and repeated-card case
study. \\
FR-4 & Interact with hosting pages. & Hosting agent, interact tool, popup
recovery, media-state checks. & Player activation KPI and popup-interference
case study. \\
FR-5 & Inspect embedded players. & Embedded agent, iframe context handling,
delayed harvest policy. & Embedded delayed-manifest case study and stream
evidence extraction metric. \\
FR-6 & Capture screenshots. & Screenshot tool, runtime events, screenshot
storage. & Screenshot review checklist and visual evidence requirements. \\
FR-7 & Extract stream or media evidence. & Harvest tool, network inspection,
manifest detection, player-state verification. & Stream evidence extraction
rate and tokenized-URL case study. \\
FR-8 & Enrich stream evidence with provider context. & Provider lookup stage,
IP/RDAP/WHOIS enrichment rows. & Provider coverage table and provider case
study. \\
FR-9 & Prepare reviewable email drafts. & Downstream email packaging stage. &
Provider/email prompt contract and responsibility appendix. \\
FR-10 & Support regression evaluation. & Dataset batches, persisted runs,
model/tool telemetry, KPI tables. & Prompt regression history and evaluation
protocol. \\
\end{longtable}

\section{Non-Functional Requirement Traceability}

\begin{longtable}{>{\raggedright\arraybackslash\bfseries}p{1.8cm}
                  >{\raggedright\arraybackslash}p{3.8cm}
                  >{\raggedright\arraybackslash}p{4.3cm}
                  >{\raggedright\arraybackslash}p{4.1cm}}
\caption{Non-functional requirements mapped to design responses.}
\label{tab:nonfunctional-traceability}\\
\toprule
\rowcolor{accent!8}
ID & Requirement & Design response & Evaluation or review signal \\
\midrule
\endfirsthead
\toprule
\rowcolor{accent!8}
ID & Requirement & Design response & Evaluation or review signal \\
\midrule
\endhead
\bottomrule
\endfoot
NFR-1 & Traceability & Runtime events, screenshots, tool calls, model usage,
and final outputs are stored separately. & Operator can reconstruct the run
path after completion. \\
NFR-2 & Maintainability & The workflow is split into orchestrator,
classification, landing, hosting, embedded, provider, and email stages. &
Agent responsibility matrix identifies stage ownership. \\
NFR-3 & Observability & Run summary, event stream, model usage, tool health,
provider rows, and screenshots are visible through the operator console. &
Dashboard telemetry and runtime evidence appendix. \\
NFR-4 & Reliability & Regression batches check old and new websites after each
prompt, model, or tool change. & Regression retention and prompt-version
tables. \\
NFR-5 & Cost control & Model usage separates input, cached input, output, and
estimated spend. & Cost per run, cache hit rate, and token summaries. \\
NFR-6 & Portability & OWC separates API, browser
tooling, workers, database, and external services. & Deployment direction and
component architecture. \\
NFR-7 & Human review & The system prepares evidence and drafts but does not
send takedown messages automatically. & Operator review flow and final
limitations. \\
\end{longtable}

\section{Evaluation Traceability}

The evaluation metrics were selected because they correspond to real evidence
risks. A metric was kept only when it helped diagnose a failure or compare two
versions of the system.

\begin{table}[!htbp]
\centering
\begin{tabularx}{\textwidth}{Q{4.0cm}Y}
\toprule
\rowcolor{accent!8}
Metric & Requirement it supports \\
\midrule
Website success rate & Confirms that a run reaches end-to-end reviewable
evidence, not only a completed status. \\
Classification accuracy & Validates the first routing decision and protects the
specialist agents from wrong tasks. \\
Host-candidate recall & Measures whether landing pages are explored deeply
enough. \\
Player activation success & Measures whether hosting and embedded agents
perform the right browser actions. \\
Stream evidence extraction rate & Confirms that the run returns useful media,
manifest, embedded, or network evidence. \\
Provider coverage & Measures whether extracted evidence can be enriched with
infrastructure context. \\
Tool success rate & Detects operational instability in the browser and provider
tools. \\
Regression retention & Checks that new changes do not break websites that
already worked. \\
Cost per working website & Links model cost to useful evidence, not just token
volume. \\
\bottomrule
\end{tabularx}
\caption{Evaluation metrics mapped to project requirements.}
\label{tab:evaluation-traceability}
\end{table}

\section{Risk Traceability}

\begin{longtable}{>{\raggedright\arraybackslash\bfseries}p{3.4cm}
                  >{\raggedright\arraybackslash}p{4.5cm}
                  >{\raggedright\arraybackslash}p{5.6cm}}
\caption{Main technical risks and mitigation choices.}
\label{tab:risk-traceability}\\
\toprule
\rowcolor{accent!8}
Risk & Impact on the project & Mitigation used in the final design \\
\midrule
\endfirsthead
\toprule
\rowcolor{accent!8}
Risk & Impact on the project & Mitigation used in the final design \\
\midrule
\endhead
\bottomrule
\endfoot
Dynamic page state & Static requests miss player and stream behavior. & Real
browser automation with inspection, interaction, screenshots, and network
harvest. \\
Popup interference & The agent may reason over an advertisement tab. & Popup
close, focus restore, and visual checkpoints after clicks. \\
Wrong route & A specialist receives the wrong task and uses the wrong tools. &
Classification stage, handoff context, and page-role definitions. \\
Premature harvesting & The stream does not exist yet when the tool runs. &
Activation guard, media-state check, and delayed harvest. \\
Short-lived stream URLs & Evidence may expire before manual replay. &
Timestamped stream records, screenshots, provider context, and trace storage. \\
Model cost growth & Repeated batches can become expensive. & Cached-input
tracking, cost per run, and early success exit. \\
Evaluation drift & A change fixes one site but breaks another. & Regression
batches and retention checks. \\
Incomplete provider data & Provider may resolve without abuse contact. &
Separate provider, IP, and abuse-contact metrics. \\
\end{longtable}

\section{Traceability Conclusion}

The traceability matrix confirms that the work is not a collection of unrelated
features. The requirements, architecture, implementation, and evaluation are
connected by the same evidence objective: make each investigation reviewable,
repeatable, and measurable. This is the reason the final report emphasizes
state, handoff, screenshots, tool timing, provider enrichment, cost, and
regression retention throughout the chapters.
